\chapter{Your First Agent}
\label{ch:first-agent}

\begin{quote}
\itshape
Three lines of meaningful code, and about forty lines of scaffolding that will
save you a fortnight later.
\end{quote}

\noindent
This chapter builds the solution every later chapter adds to, gets an agent
answering questions against a model running on your own machine, and then does
the thing that makes the rest of the book cheap to follow: proves that the same
agent runs against several different providers without touching the agent code.

Along the way it introduces conversation threads, which look like a minor
convenience and are in fact the mechanism that makes Chapters~\ref{ch:memory}
and~\ref{ch:workflows-advanced} possible.

\begin{verifybox}
The listings in this chapter are transcribed from Microsoft's official
documentation and the \code{dotnet/samples/01-get-started} directory of the
framework repository. They have not been compiled by the author against
\mafcore{}. Two of them --- flagged individually --- involve API names that
changed between preview and GA and are still inconsistent across published
material. Compile before you copy.
\end{verifybox}

% ============================================================
\section{What you will have at the end}

\begin{itemize}[leftmargin=1.4em]
\item A solution with Central Package Management, one pinned version set, and a
      folder per chapter.
\item A model serving locally, with no cloud subscription and no API key.
\item An agent that answers in one shot and an agent that streams.
\item A conversation that survives a process restart.
\item A provider-swap harness, and the beginning of a measurement table that
      will be worth more to you than any of the code.
\end{itemize}

% ============================================================
\section{The solution scaffold}
\label{sec:scaffold}
\index{Central Package Management}
\index{project templates}

\subsection{Start from a template, or by hand?}
\label{sec:templates}

Before typing anything, know that a template package exists. Three, in fact, and
they are easy to miss because none of them is installed by default:

\begin{center}
\small
\begin{tabularx}{\linewidth}{@{}lX@{}}
\toprule
\textbf{Package} & \textbf{Provides} \\
\midrule
\pkg{Microsoft.Agents.AI.ProjectTemplates} & An AI Agent Web API project --- an ASP.NET Core host with an agent already wired up. \\
\pkg{Microsoft.Extensions.AI.Templates} & AI chat web applications, with vector storage and ingestion for retrieval. \\
\pkg{Microsoft.McpServer.ProjectTemplates} & An MCP server project. Chapter~\ref{ch:tools} writes one of these by hand. \\
\bottomrule
\end{tabularx}
\end{center}

\begin{shellcmd}
dotnet new install Microsoft.Agents.AI.ProjectTemplates::1.3.0-preview.1.26251.3
dotnet new list agent
\end{shellcmd}

\begin{verifybox}
The short names the templates register under have not been verified by the
author. Run \code{dotnet new list} after installing and write down what you
actually get --- the package version above is current as of \mafdate{} and the
template names may well have changed by the time you read this.
\end{verifybox}

The agent web API template asks which provider you want and can wire up Azure
OpenAI, OpenAI, GitHub Models or Ollama, with .NET Aspire integration available.

\begin{note}
\textbf{GitHub Models is worth knowing about} if your machine cannot comfortably
serve a local model. It offers hosted models against a GitHub personal access
token with a \code{models} scope, which makes it a third option alongside ``pay
for Azure'' and ``run it on your own GPU''. The samples in this book default to
Ollama, but the provider factory in \S\ref{sec:swap} accommodates any of them.
\end{note}

\needscreenshot{vs-new-project-agent-template}%
{The agent templates in the New Project dialog}%
{\vsver{}: \emph{File $\rightarrow$ New $\rightarrow$ Project}, then type
\texttt{agent} into the template search box, \textbf{after} installing
\texttt{Microsoft.Agents.AI.ProjectTemplates}. Capture the dialog showing the
filtered template list with the Agent Framework entries visible and one selected
so its description panel is readable on the right. If the templates do not appear
in the dialog after a CLI install, restart Visual Studio first --- and if they
still do not appear, capture the \texttt{dotnet new list} terminal output instead
and name the file the same thing.}

\subsection{Why this book scaffolds by hand anyway}

Three reasons, and the third is the one that matters.

\textbf{The templates produce one project; you need a solution.} This book
accumulates roughly a dozen projects across eight chapters, sharing one version
set. A template that scaffolds a single self-contained web API is solving a
different problem.

\textbf{The templates are preview and lightly adopted.} At the time of writing
the agent template package had a few thousand downloads against several million
for the core package. That ratio means the template will drift relative to the
framework, and drift in a scaffold is expensive because you only notice it when
something else breaks.

\textbf{The scaffolding is the lesson.} What follows --- Central Package
Management across release trains --- is precisely the thing a template hides,
and precisely the thing that will cost you an afternoon in
Chapter~\ref{ch:hosting} when you add a package from a different train. Typing it
once means you understand the failure when it arrives.

\begin{exercisebox}
Install the templates anyway and generate one agent web API project into a
scratch directory. Do not delete it immediately. Read the generated
\code{Program.cs} and \code{.csproj}, and answer two questions: does the template
pin package versions, and does it use Central Package Management? Then compare
its agent construction to Listing~\ref{lst:seam}.

This takes ten minutes and tells you what Microsoft considers a default project
shape --- which is useful context whether or not you adopt it.
\end{exercisebox}

\subsection{Creating the solution}

Start from the command line; the IDE catches up afterwards.

\begin{shellcmd}
mkdir maf-book-samples && cd maf-book-samples
dotnet new sln -n MafBookSamples
dotnet new gitignore

mkdir -p src/01-agents
cd src/01-agents
dotnet new console -o HelloAgent
cd ../..
dotnet sln add src/01-agents/HelloAgent/HelloAgent.csproj
\end{shellcmd}

\subsection{Pinning versions in one place}

Agent Framework ships as roughly three dozen packages on several release trains,
and by Chapter~\ref{ch:hosting} you will have a dozen projects referencing
overlapping sets of them. Central Package Management solves this before it
becomes a problem. Create \code{Directory.Packages.props} at the solution root:

\begin{xmlcode}[caption={\code{Directory.Packages.props} --- the only place versions are written}, label={lst:cpm}]
<Project>
  <PropertyGroup>
    <ManagePackageVersionsCentrally>true</ManagePackageVersionsCentrally>
    <CentralPackageTransitivePinningEnabled>true</CentralPackageTransitivePinningEnabled>
  </PropertyGroup>

  <ItemGroup Label="Agent Framework">
    <PackageVersion Include="Microsoft.Agents.AI" Version="1.10.0" />
    <PackageVersion Include="Microsoft.Agents.AI.Abstractions" Version="1.10.0" />
    <PackageVersion Include="Microsoft.Agents.AI.OpenAI" Version="1.10.0" />
    <PackageVersion Include="Microsoft.Agents.AI.Workflows" Version="1.10.0" />
  </ItemGroup>

  <ItemGroup Label="Providers">
    <PackageVersion Include="Azure.AI.OpenAI" Version="2.1.0" />
    <PackageVersion Include="Azure.Identity" Version="1.13.1" />
    <PackageVersion Include="Microsoft.Extensions.AI" Version="9.4.0" />
    <PackageVersion Include="Microsoft.Extensions.AI.Ollama" Version="9.4.0-preview.1" />
  </ItemGroup>

  <ItemGroup Label="Hosting and configuration">
    <PackageVersion Include="Microsoft.Extensions.Configuration" Version="9.0.0" />
    <PackageVersion Include="Microsoft.Extensions.Configuration.UserSecrets" Version="9.0.0" />
    <PackageVersion Include="Microsoft.Extensions.DependencyInjection" Version="9.0.0" />
  </ItemGroup>
</Project>
\end{xmlcode}

\begin{versionbox}
Replace the version numbers in Listing~\ref{lst:cpm} with whatever is current
when you read this. The four core packages above share a version because they
ship on the same train. The moment you add a package from the extensions train
--- the harness in Chapter~\ref{ch:harness}, durable execution in
Chapter~\ref{ch:hosting} --- you will be pinning a different number in the same
file, and that is correct rather than a mistake.
Appendix~\ref{app:packages} maps every package to its train.
\end{versionbox}

Individual projects then reference packages without versions:

\begin{xmlcode}[caption={\code{HelloAgent.csproj} --- no version numbers anywhere}, label={lst:csproj}]
<Project Sdk="Microsoft.NET.Sdk">
  <PropertyGroup>
    <OutputType>Exe</OutputType>
    <TargetFramework>net10.0</TargetFramework>
    <ImplicitUsings>enable</ImplicitUsings>
    <Nullable>enable</Nullable>
  </PropertyGroup>

  <ItemGroup>
    <PackageReference Include="Microsoft.Agents.AI" />
    <PackageReference Include="Microsoft.Extensions.AI.Ollama" />
  </ItemGroup>
</Project>
\end{xmlcode}

\needscreenshot{vs-cpm-directory-packages}%
{Central Package Management: one file, every version}%
{\vsver{} with \texttt{Directory.Packages.props} open in the editor and Solution
Explorer visible on the right showing the solution root with
\texttt{Directory.Packages.props}, \texttt{Directory.Build.props} and the
\texttt{src} folder. The point of the capture is the relationship between the two
panes, so include both. Scroll so that all three \texttt{ItemGroup} labels are
visible at once.}

\begin{warning}
\index{prerelease packages}
\pkg{Microsoft.Extensions.AI.Ollama} is a \emph{prerelease} package even though
Agent Framework itself is GA --- the two live in different package families with
different release cadences. If \code{dotnet restore} tells you the package does
not exist, you have not enabled prerelease resolution. In Visual Studio this is
the \emph{Include prerelease} checkbox; on the command line it is
\code{-{}-prerelease} on \code{dotnet add package}. This trips up almost
everyone once.
\end{warning}

\needscreenshot{vs-nuget-prerelease-ollama}%
{The prerelease checkbox, which is not optional here}%
{\vsver{} \emph{Manage NuGet Packages} on the \texttt{HelloAgent} project,
\emph{Browse} tab, search \texttt{Microsoft.Extensions.AI.Ollama}. Capture with
the \emph{Include prerelease} checkbox \textbf{ticked} and the package visible in
the results, with the version dropdown on the right showing a \texttt{-preview}
version selected. If convenient, take a second capture with the box unticked and
no results, so the two can be shown side by side --- name that one
\texttt{vs-nuget-prerelease-ollama-unticked.png}.}

\subsection{Shared build settings}

One more file at the root, \code{Directory.Build.props}, so that every project
in every chapter shares the same language settings and you never adjust them
project by project:

\begin{xmlcode}[caption={\code{Directory.Build.props}}, label={lst:buildprops}]
<Project>
  <PropertyGroup>
    <TargetFramework>net10.0</TargetFramework>
    <LangVersion>latest</LangVersion>
    <Nullable>enable</Nullable>
    <ImplicitUsings>enable</ImplicitUsings>
    <TreatWarningsAsErrors>true</TreatWarningsAsErrors>
    <InvariantGlobalization>true</InvariantGlobalization>
  </PropertyGroup>
</Project>
\end{xmlcode}

The finished structure is the one shown in
Figure~\ref{fig:vs-solution-explorer-scaffold} in the previous chapter.

% ============================================================
\section{A model on your own machine}
\label{sec:local-model}
\index{Ollama}

Everything in Parts~I and~II runs locally. This is not a gesture toward readers
without an Azure subscription --- it is a practical constraint. You are about to
run the same prompt several hundred times while learning, and the difference
between that costing nothing and costing real money determines whether you
actually experiment or merely read.

Install Ollama, then pull a model:

\begin{shellcmd}
ollama pull llama3.2
ollama list
curl http://localhost:11434/api/tags
\end{shellcmd}

Ollama serves on port \code{11434} by default and needs no authentication.

\begin{warning}
\index{function calling!model support}
\textbf{Not every local model supports function calling.} This matters from
Chapter~\ref{ch:tools} onward, and it produces a confusing failure --- the agent
answers, politely, without ever calling your tool, and nothing errors. Stick to
models documented as supporting tool use; \code{llama3.2} and \code{qwen3:4b} are
the two Microsoft's own provider documentation names. Verify tool calling works
before you conclude your tool registration is broken.
\end{warning}

\begin{note}
Model choice affects the quality of every experiment in this book, and the
smaller local models will underperform a frontier hosted model noticeably on
multi-step reasoning. That is fine for learning mechanics and misleading for
evaluating capability. Chapter~\ref{ch:evaluation} returns to this properly.
Where a result genuinely depends on model strength, the text says so.
\end{note}

% ============================================================
\section{Constructing an agent}
\label{sec:construction}

There are three construction forms in circulation. They produce the same thing.
Knowing why there are three saves you from the confusion of finding all of them
in different articles.

\subsection{The fluent one-liner}

The form used throughout Microsoft's Azure OpenAI samples chains authentication
straight into an agent:

\begin{csharp}[caption={The canonical Azure OpenAI agent}, label={lst:azure-agent}]
using Azure.AI.OpenAI;
using Azure.Identity;
using Microsoft.Agents.AI;

var endpoint = Environment.GetEnvironmentVariable("AZURE_OPENAI_ENDPOINT")
    ?? throw new InvalidOperationException("AZURE_OPENAI_ENDPOINT is not set.");
var deploymentName = Environment.GetEnvironmentVariable("AZURE_OPENAI_DEPLOYMENT_NAME")
    ?? "gpt-4o-mini";

AIAgent agent = new AzureOpenAIClient(
        new Uri(endpoint),
        new DefaultAzureCredential())
    .GetChatClient(deploymentName)
    .AsAIAgent(
        instructions: "You are good at telling jokes.",
        name: "Joker");

Console.WriteLine(await agent.RunAsync("Tell me a joke about a pirate."));
\end{csharp}

The chain reads: authentication $\rightarrow$ chat client $\rightarrow$ agent. It
is compact, and it is also the least instructive of the three, because it hides
the seam the whole framework is built on.

\subsection{The seam: \texorpdfstring{\texttt{IChatClient}}{IChatClient}}
\label{sec:seam}
\index{IChatClient@\texttt{IChatClient}}

Split the chain in two and the seam becomes visible:

\begin{csharp}[caption={Provider-specific above the line, provider-agnostic below}, label={lst:seam}]
using Microsoft.Agents.AI;
using Microsoft.Extensions.AI;

// ---- everything above this line is provider-specific ----
IChatClient chatClient = new OllamaChatClient(
    new Uri("http://localhost:11434"),
    modelId: "llama3.2");
// ---- everything below this line is provider-agnostic ----

AIAgent agent = chatClient.AsAIAgent(
    instructions: "You are a helpful assistant running locally via Ollama.",
    name: "LocalAgent");

Console.WriteLine(await agent.RunAsync("What is the largest city in France?"));
\end{csharp}

This is the form the book uses, and Listing~\ref{lst:seam} is worth staring at
for a moment. \api{IChatClient} comes from \pkg{Microsoft.Extensions.AI} --- it
is not an Agent Framework type at all. Agent Framework consumes it. That is the
entire mechanism behind the provider swap in \S\ref{sec:swap}: there is no
provider abstraction inside Agent Framework, because the layer below already
supplied one.

\begin{note}
There are two ways to reach Ollama from .NET, and both appear in published
samples. \pkg{Microsoft.Extensions.AI.Ollama} provides \api{OllamaChatClient},
used above and in Microsoft's provider documentation. The community
\pkg{OllamaSharp} package provides \api{OllamaApiClient}, which also implements
\api{IChatClient} and therefore works identically. Either is fine. This book uses
the first because it is first-party.
\end{note}

\subsection{The concrete type: \texorpdfstring{\texttt{ChatClientAgent}}{ChatClientAgent}}
\index{ChatClientAgent@\texttt{ChatClientAgent}}

\api{AsAIAgent} is an extension method. What it constructs is a
\api{ChatClientAgent}, which you can also new up directly:

\begin{csharp}[caption={The same agent, constructed explicitly}, label={lst:chatclientagent}]
using OllamaApiClient chatClient = new(new Uri("http://localhost:11434"), "phi3");

AIAgent agent = new ChatClientAgent(
    chatClient,
    instructions: "You are good at telling jokes.",
    name: "Joker");
\end{csharp}

\api{ChatClientAgent} is the general-purpose agent type: it takes any
\api{IChatClient} and adds instruction handling, conversation history, and the
tool-calling loop. Reach for the explicit constructor when you need the options
object --- for a custom message store, for instance, which is how
Chapter~\ref{ch:memory} persists conversations to PostgreSQL.

\begin{csharp}[caption={The options form, used from Chapter~\ref{ch:memory} onward}, label={lst:agent-options}]
var agent = chatClient.CreateAIAgent(new ChatClientAgentOptions
{
    Name = "Assistant",
    ChatOptions = new()
    {
        Instructions = "You are a helpful assistant."
    },
    // Plug in your own history store; the framework calls it to load and save.
    ChatMessageStoreFactory = ctx => new FileChatMessageStore(ctx.SerializedState)
});
\end{csharp}

\begin{verifybox}
\textbf{Two names for one operation.} You will find both \api{AsAIAgent} and
\api{CreateAIAgent} in current published material, sometimes in documentation
written weeks apart. This is residue from the preview-to-GA transition rather
than two distinct APIs. Check which one your pinned version exposes ---
IntelliSense on the \api{IChatClient} extension methods will settle it in five
seconds --- and use that one consistently. Listing~\ref{lst:agent-options} in
particular should be verified before use.
\end{verifybox}

\subsection{On instructions and names}

The \code{instructions} argument is the agent's system prompt, fixed at
construction. Two consequences follow, and both matter later.

First, \textbf{instructions are not dynamic}. If you need them to vary per
request, construct a new agent per instruction set, or inject messages at the
session level. Since an agent is cheap --- a thin decoration over a chat client
--- constructing several is the expected move rather than a workaround.

Second, \textbf{the name is functional, not decorative}. In handoff orchestration
(Chapter~\ref{ch:orchestration}) the framework generates transfer tools from
agent names, which means the string you type becomes part of the model's decision
surface. \code{Billing} routes better than \code{Agent2}. Get into the habit now.

% ============================================================
\section{Running: blocking and streaming}
\label{sec:running}

Two methods, one distinction.

\begin{csharp}[caption={Both response modes}, label={lst:streaming}]
// Blocking: waits for the complete response.
Console.WriteLine("=== Non-Streaming ===");
Console.WriteLine(await agent.RunAsync("Tell me a joke about a pirate."));

// Streaming: yields updates as the model produces them.
Console.WriteLine("=== Streaming ===");
await foreach (var update in agent.RunStreamingAsync("Tell me a joke about a pirate."))
{
    Console.Write(update);   // no newline, so output accumulates in place
}
Console.WriteLine();
\end{csharp}

\api{RunAsync} returns the complete response. \api{RunStreamingAsync} returns an
async sequence of updates.

The choice is not stylistic. Use streaming wherever a human is waiting, because
time-to-first-token is what the user actually experiences; use blocking for
background processing, or wherever the next step needs the whole response before
it can begin --- which, notably, includes most workflow executors in Part~II.

\begin{note}
Writing \code{Console.Write(update)} works because the update type has a sensible
\api{ToString}. In a real UI you will want the structured content rather than the
string, and the shape of an update differs slightly between providers. This is
one of the seams \S\ref{sec:swap-limits} examines.
\end{note}

% ============================================================
\section{Conversation threads}
\label{sec:threads}
\index{AgentThread@\texttt{AgentThread}}

Every example so far has been amnesiac. Ask a follow-up question and the agent
has no idea what you are referring to, because each \api{RunAsync} call without a
thread creates a throwaway one that lives only for the duration of that call.

\api{AgentThread} is the fix, and it is a more interesting type than its name
suggests. It holds the conversation history \emph{plus} any other state the agent
must preserve across interactions --- including memories, and including
references to state that lives somewhere else entirely.

\subsection{Creating and using a thread}

\begin{csharp}[caption={A thread makes the conversation continuous}, label={lst:thread-basic}]
AgentThread thread = agent.GetNewThread();

Console.WriteLine(await agent.RunAsync("My name is Konrad.", thread));
Console.WriteLine(await agent.RunAsync("What is my name?", thread));
\end{csharp}

Threads are cheap and independent. One agent can serve many concurrent
conversations:

\begin{csharp}[caption={One agent, two independent conversations}, label={lst:two-threads}]
AgentThread thread1 = agent.GetNewThread();
AgentThread thread2 = agent.GetNewThread();

Console.WriteLine(await agent.RunAsync("What weights should I start with?", thread1));
Console.WriteLine(await agent.RunAsync("How long will 5x5 take me?", thread2));
\end{csharp}

That separation is what makes a single registered agent instance viable behind a
multi-user web application, which is exactly how Chapter~\ref{ch:hosting} uses it.

\subsection{Serialising and resuming}

\begin{csharp}[caption={A conversation that survives the process}, label={lst:thread-serialize}]
AgentThread thread = agent.GetNewThread();
var response = await agent.RunAsync("Hello, how are you?", thread);

// Serialise for storage.
JsonElement serializedThread = thread.Serialize();
var json = JsonSerializer.Serialize(serializedThread,
    new JsonSerializerOptions { WriteIndented = true });
await File.WriteAllTextAsync("thread.json", json);

// ... process restarts ...

var reloaded = JsonSerializer.Deserialize<JsonElement>(
    await File.ReadAllTextAsync("thread.json"));
AgentThread resumedThread = agent.DeserializeThread(reloaded);

response = await agent.RunAsync("What did I just ask you?", resumedThread);
\end{csharp}

\begin{verifybox}
\textbf{Thread, or session?} This is the most consequential naming instability in
the framework at the time of writing, and it affects every listing in this
section.

Older material --- including much of what is still published --- uses
\api{AgentThread}, \api{GetNewThread}, \api{DeserializeThread} and
\api{AgentRunResponse}. Microsoft's most recent documentation pages instead use
\api{AgentSession}, \api{CreateSessionAsync} and \api{AgentResponse}:

\begin{quote}
\code{AgentSession session = await agent.CreateSessionAsync();}\\
\code{AgentResponse response = await agent.RunAsync("...", session);}
\end{quote}

The framework overview now describes ``an agent session for state management,''
which suggests the session naming is the destination rather than a parallel API.
Some intermediate versions also expose \api{GetNewThreadAsync} and
\api{DeserializeThreadAsync}.

Check what your pinned version actually exposes before writing anything on top of
it, and pick one vocabulary for your whole codebase. Where this book says
\emph{thread}, read \emph{session} if that is what your SDK gives you --- the
concept is identical and everything else in this section holds.
\end{verifybox}

Three things about that serialised blob are worth knowing before
Chapter~\ref{ch:memory} builds a real store on top of it.

\textbf{Its contents depend on the provider.} Where history is held in memory,
the serialised thread contains the messages themselves. Where the provider stores
conversations server-side --- Foundry Agents, OpenAI Responses --- the serialised
thread contains an \emph{identifier} pointing at the remote conversation. Same
API, radically different persistence characteristics, and different
data-residency implications.

\textbf{Service-side threads are your responsibility to clean up.} If the
provider created something remote, deleting it is not automatic. A long-running
service that creates threads and never removes them accumulates state you are
paying for.

\begin{warning}
\textbf{You must serialise and deserialise with the same agent instance.}
Deserialising a thread on a differently configured agent is not supported, and
the failure will not necessarily be loud. In practice this means agent
construction must be deterministic --- registered once in DI, with configuration
that does not drift between the process that saved and the process that loads.
\end{warning}

\subsection{Looking inside a thread}

Before moving on, put a breakpoint after the second \api{RunAsync} and inspect
the thread. It is the fastest way to build an accurate mental model of what the
framework actually sends to the model on each turn --- and, later, of what
context compaction removes.

\needscreenshot{vs-debugger-agentthread-watch}%
{Inspecting an \texttt{AgentThread} mid-conversation}%
{\vsver{} stopped at a breakpoint on the line after the second
\texttt{RunAsync} in Listing~\ref{lst:thread-basic}. Capture the editor with the
breakpoint and the yellow current-line marker visible, plus the \emph{Locals} or
\emph{Watch} window with the \texttt{thread} variable expanded far enough to show
the accumulated message collection --- ideally showing both the user messages and
the assistant replies. Widen the Locals pane so values are not truncated.}

% ============================================================
\section{Swapping providers}
\label{sec:swap}

Because the provider abstraction lives below Agent Framework
(\S\ref{sec:seam}), swapping providers means changing which \api{IChatClient} you
construct and nothing else.

\begin{csharp}[caption={A provider factory --- the only provider-aware code in the solution}, label={lst:factory}]
using Azure.AI.OpenAI;
using Azure.Identity;
using Microsoft.Extensions.AI;

static IChatClient CreateChatClient(IConfiguration config) =>
    config["Provider"] switch
    {
        "ollama" => new OllamaChatClient(
            new Uri(config["Ollama:Endpoint"] ?? "http://localhost:11434"),
            modelId: config["Ollama:Model"] ?? "llama3.2"),

        "azureopenai" => new AzureOpenAIClient(
                new Uri(config["AzureOpenAI:Endpoint"]!),
                new DefaultAzureCredential())
            .GetChatClient(config["AzureOpenAI:Deployment"]!)
            .AsIChatClient(),

        var other => throw new NotSupportedException($"Unknown provider '{other}'.")
    };
\end{csharp}

Note \api{AsIChatClient} on the Azure branch. It comes from
\pkg{Microsoft.Extensions.AI.OpenAI} and converts the provider-specific chat
client into the standard interface. In Listing~\ref{lst:azure-agent} it was
implicit, because \api{AsAIAgent} was applied directly to the provider client.
Here it is explicit, because the factory must return a common type.

Everything downstream is unchanged:

\begin{csharp}[caption={Agent construction, now provider-blind}, label={lst:provider-blind}]
IChatClient chatClient = CreateChatClient(config);

AIAgent agent = chatClient.AsAIAgent(
    instructions: "You are a helpful assistant.",
    name: "Assistant");
\end{csharp}

\begin{note}
This pattern is worth more than the convenience it looks like. It means your full
test suite can run against a local model in CI at zero marginal cost, while
production runs against a hosted one. Chapter~\ref{ch:evaluation} builds directly
on that.
\end{note}

\subsection{Secrets, not environment variables}

Listing~\ref{lst:azure-agent} reads configuration from environment variables
because that is how Microsoft's samples do it. Do not carry that habit into
anything you commit.

\begin{shellcmd}
dotnet user-secrets init
dotnet user-secrets set "AzureOpenAI:Endpoint" "https://your-resource.openai.azure.com/"
dotnet user-secrets set "AzureOpenAI:Deployment" "gpt-4o-mini"
\end{shellcmd}

Prefer \api{DefaultAzureCredential} over key-based authentication wherever the
provider supports it --- there is then no secret to leak.

\needscreenshot{vs-user-secrets-manage}%
{User secrets, so nothing lands in source control}%
{\vsver{}: right-click the \texttt{HelloAgent} project, \emph{Manage User
Secrets}, which opens \texttt{secrets.json}. Capture the editor showing the file
with the \texttt{AzureOpenAI} section present. \textbf{Redact the endpoint host
and any deployment name that identifies a real resource} before capturing, or use
a throwaway resource name. If you want to show how to get there, take the context
menu as a second capture named \texttt{vs-user-secrets-menu.png}.}

\subsection{Where the one-line swap stops being one line}
\label{sec:swap-limits}

The marketing claim is that swapping providers is a one-line change. At the
constructor, it is true. At the seams, it is not --- and finding out where in a
controlled experiment now is far better than finding out in production later.

Four places the abstraction leaks:

\begin{description}[leftmargin=0pt, style=nextline]
\item[Tool-calling fidelity.] Providers differ in how reliably they select the
right tool and how well they handle complex parameter schemas. Small local models
differ from frontier models here by a wide margin. This is the biggest practical
difference, and Chapter~\ref{ch:tools} measures it directly.

\item[Streaming granularity.] Some providers emit token-level updates; others
emit larger chunks. Code that assumes fine-grained streaming produces a visibly
worse experience against a provider that batches.

\item[Token accounting.] What is reported, and when, differs. If you are
computing cost from usage data --- and by Chapter~\ref{ch:observability} you will
be --- verify per provider rather than assuming a shared shape.

\item[Thread persistence semantics.] As \S\ref{sec:threads} covered: in-memory
messages for some providers, a remote identifier for others. This changes your
storage design, not merely your storage size.
\end{description}

% ============================================================
\section{Measuring the swap}
\label{sec:measure}

This section is the reason the chapter exists in this form, and it is the one
piece of work here that is not available anywhere else.

Build a small harness that runs a fixed prompt set against each configured
provider and records, per provider: time to first token, total wall time, tokens
reported in and out, and --- once Chapter~\ref{ch:tools} adds tools --- tool
selection accuracy over repeated runs.

\begin{exercisebox}
Run a fixed set of ten prompts (five conversational, five requiring a specific
structured answer) against at least three providers, twenty times each. Record:

\begin{itemize}[leftmargin=1.4em]
\item median and p95 time to first token
\item median total wall time
\item reported prompt and completion tokens
\item variance in output length for identical inputs
\end{itemize}

Then answer the question the numbers exist to answer: \emph{for which of your own
use cases would the cheapest provider be sufficient?} Most teams never ask this,
and default to the most expensive option available.
\end{exercisebox}

Appendix~\ref{app:providers} carries the author's measurements. Yours will differ
--- hardware, model versions and region all move the numbers --- which is
precisely why running the harness yourself is the exercise rather than reading a
table.

% ============================================================
\section{Summary}

An agent is an \api{IChatClient} plus instructions, a name, and (from
Chapter~\ref{ch:tools}) tools. Three construction forms exist and produce the same
\api{ChatClientAgent}; the two-step form via an explicit \api{IChatClient} is the
one to standardise on, because it makes the provider seam visible.

\api{RunAsync} blocks and returns the whole response; \api{RunStreamingAsync}
yields updates. Use the second wherever a human is waiting.

\api{AgentThread} carries conversation history and any other state that must
survive between runs. It serialises to a \api{JsonElement}, must be resumed by an
identically configured agent instance, and contains either the messages
themselves or a reference to server-side state depending on the provider --- a
distinction Chapter~\ref{ch:memory} builds on.

Provider swapping is genuinely a one-line change at the constructor and genuinely
not at the seams. The four leaks are tool-calling fidelity, streaming
granularity, token accounting, and thread persistence semantics.

\begin{exercisebox}
Before Chapter~\ref{ch:tools}:

\begin{enumerate}[leftmargin=1.4em]
\item Get \code{HelloAgent} answering against a local model with no cloud account.
\item Add a second provider behind the factory in Listing~\ref{lst:factory} and
      verify the agent code did not change.
\item Hold a four-turn conversation, serialise the thread to disk, restart the
      process, resume it, and confirm the agent still knows what you told it in
      turn one.
\item Open the serialised \code{thread.json} in an editor. Write down, in your own
      words, what is in it --- and predict what would change if the provider stored
      conversations server-side.
\end{enumerate}
\end{exercisebox}
